% Part 1-3 section generated by problem_set_02/generate_part1_3_writeup.py
% This file is intended to be % Part 1-3 section generated by problem_set_02/generate_part1_3_writeup.py
% This file is intended to be % Part 1-3 section generated by problem_set_02/generate_part1_3_writeup.py
% This file is intended to be % Part 1-3 section generated by problem_set_02/generate_part1_3_writeup.py
% This file is intended to be \input{problem_set_02/part1_3_writeup.tex} into a larger document.

\section*{Problem Set 02: Parts 1--3}

\subsection*{Part 1: Preferences and Priorities Summary}
The market instance uses 18 students, 3 schools, and school capacity 6 (total capacity 18), with strict random preferences and strict random school priorities. The one-market realization is generated with \texttt{random.seed(42)}. Detailed preference and priority lists are documented in \texttt{problem\_set\_02/market\_setup\_output.md}.

\subsection*{Part 2: One-Market Matching Table (Seed = 42)}
The table below reports student assignments under Deferred Acceptance (DA), Immediate Acceptance (IA/Boston), and Top Trading Cycles (TTC), all run on the same Part 1 market instance.

\begin{center}
\begin{tabular}{lccc}
\hline
Student & DA & IA & TTC \\
\hline
i1 & s2 & s2 & s2 \\
i2 & s3 & s3 & s3 \\
i3 & s3 & s3 & s3 \\
i4 & s3 & s3 & s3 \\
i5 & s1 & s1 & s1 \\
i6 & s1 & s1 & s1 \\
i7 & s3 & s3 & s3 \\
i8 & s2 & s2 & s2 \\
i9 & s2 & s2 & s2 \\
i10 & s2 & s2 & s2 \\
i11 & s1 & s1 & s1 \\
i12 & s3 & s3 & s3 \\
i13 & s1 & s1 & s1 \\
i14 & s2 & s2 & s2 \\
i15 & s1 & s1 & s1 \\
i16 & s1 & s1 & s1 \\
i17 & s2 & s2 & s2 \\
i18 & s3 & s3 & s3 \\
\hline
\end{tabular}
\end{center}

\paragraph{One-Market Average Rank Check}
\begin{center}
\begin{tabular}{lc}
\hline
Mechanism & Average Rank in this one market \\
\hline
DA & 1.2222 \\
IA & 1.2222 \\
TTC & 1.2222 \\
\hline
\end{tabular}
\end{center}

\subsection*{Part 3: Simulation Results (N = 1000)}
Each simulation draw generates a fresh random market, then computes DA/IA/TTC outcomes and student average rank. The simulation uses fixed seed \texttt{seed=123} for reproducibility.

\begin{center}
\begin{tabular}{lc}
\hline
Mechanism & Average Rank (lower is better) \\
\hline
DA & 1.2086 \\
IA & 1.1837 \\
TTC & 1.1934 \\
\hline
\end{tabular}
\end{center}

\paragraph{Interpretation}
Across 1000 simulated markets, IA has the lowest average rank (1.1837), so it performs best by this preference-based metric. TTC is close behind at 1.1934 (a gap of 0.0097), indicating similar overall assignment quality in expectation. DA is higher at 1.2086, which is 0.0249 worse than IA on average. Because lower rank means students are assigned closer to their top choices, these differences summarize welfare trade-offs in a compact, reproducible way.

\paragraph{Reproducibility}
From the repository root, run \texttt{python problem\_set\_02/generate\_part1\_3\_writeup.py} to regenerate both Markdown and LaTeX outputs.
 into a larger document.

\section*{Problem Set 02: Parts 1--3}

\subsection*{Part 1: Preferences and Priorities Summary}
The market instance uses 18 students, 3 schools, and school capacity 6 (total capacity 18), with strict random preferences and strict random school priorities. The one-market realization is generated with \texttt{random.seed(42)}. Detailed preference and priority lists are documented in \texttt{problem\_set\_02/market\_setup\_output.md}.

\subsection*{Part 2: One-Market Matching Table (Seed = 42)}
The table below reports student assignments under Deferred Acceptance (DA), Immediate Acceptance (IA/Boston), and Top Trading Cycles (TTC), all run on the same Part 1 market instance.

\begin{center}
\begin{tabular}{lccc}
\hline
Student & DA & IA & TTC \\
\hline
i1 & s2 & s2 & s2 \\
i2 & s3 & s3 & s3 \\
i3 & s3 & s3 & s3 \\
i4 & s3 & s3 & s3 \\
i5 & s1 & s1 & s1 \\
i6 & s1 & s1 & s1 \\
i7 & s3 & s3 & s3 \\
i8 & s2 & s2 & s2 \\
i9 & s2 & s2 & s2 \\
i10 & s2 & s2 & s2 \\
i11 & s1 & s1 & s1 \\
i12 & s3 & s3 & s3 \\
i13 & s1 & s1 & s1 \\
i14 & s2 & s2 & s2 \\
i15 & s1 & s1 & s1 \\
i16 & s1 & s1 & s1 \\
i17 & s2 & s2 & s2 \\
i18 & s3 & s3 & s3 \\
\hline
\end{tabular}
\end{center}

\paragraph{One-Market Average Rank Check}
\begin{center}
\begin{tabular}{lc}
\hline
Mechanism & Average Rank in this one market \\
\hline
DA & 1.2222 \\
IA & 1.2222 \\
TTC & 1.2222 \\
\hline
\end{tabular}
\end{center}

\subsection*{Part 3: Simulation Results (N = 1000)}
Each simulation draw generates a fresh random market, then computes DA/IA/TTC outcomes and student average rank. The simulation uses fixed seed \texttt{seed=123} for reproducibility.

\begin{center}
\begin{tabular}{lc}
\hline
Mechanism & Average Rank (lower is better) \\
\hline
DA & 1.2086 \\
IA & 1.1837 \\
TTC & 1.1934 \\
\hline
\end{tabular}
\end{center}

\paragraph{Interpretation}
Across 1000 simulated markets, IA has the lowest average rank (1.1837), so it performs best by this preference-based metric. TTC is close behind at 1.1934 (a gap of 0.0097), indicating similar overall assignment quality in expectation. DA is higher at 1.2086, which is 0.0249 worse than IA on average. Because lower rank means students are assigned closer to their top choices, these differences summarize welfare trade-offs in a compact, reproducible way.

\paragraph{Reproducibility}
From the repository root, run \texttt{python problem\_set\_02/generate\_part1\_3\_writeup.py} to regenerate both Markdown and LaTeX outputs.
 into a larger document.

\section*{Problem Set 02: Parts 1--3}

\subsection*{Part 1: Preferences and Priorities Summary}
The market instance uses 18 students, 3 schools, and school capacity 6 (total capacity 18), with strict random preferences and strict random school priorities. The one-market realization is generated with \texttt{random.seed(42)}. Detailed preference and priority lists are documented in \texttt{problem\_set\_02/market\_setup\_output.md}.

\subsection*{Part 2: One-Market Matching Table (Seed = 42)}
The table below reports student assignments under Deferred Acceptance (DA), Immediate Acceptance (IA/Boston), and Top Trading Cycles (TTC), all run on the same Part 1 market instance.

\begin{center}
\begin{tabular}{lccc}
\hline
Student & DA & IA & TTC \\
\hline
i1 & s2 & s2 & s2 \\
i2 & s3 & s3 & s3 \\
i3 & s3 & s3 & s3 \\
i4 & s3 & s3 & s3 \\
i5 & s1 & s1 & s1 \\
i6 & s1 & s1 & s1 \\
i7 & s3 & s3 & s3 \\
i8 & s2 & s2 & s2 \\
i9 & s2 & s2 & s2 \\
i10 & s2 & s2 & s2 \\
i11 & s1 & s1 & s1 \\
i12 & s3 & s3 & s3 \\
i13 & s1 & s1 & s1 \\
i14 & s2 & s2 & s2 \\
i15 & s1 & s1 & s1 \\
i16 & s1 & s1 & s1 \\
i17 & s2 & s2 & s2 \\
i18 & s3 & s3 & s3 \\
\hline
\end{tabular}
\end{center}

\paragraph{One-Market Average Rank Check}
\begin{center}
\begin{tabular}{lc}
\hline
Mechanism & Average Rank in this one market \\
\hline
DA & 1.2222 \\
IA & 1.2222 \\
TTC & 1.2222 \\
\hline
\end{tabular}
\end{center}

\subsection*{Part 3: Simulation Results (N = 1000)}
Each simulation draw generates a fresh random market, then computes DA/IA/TTC outcomes and student average rank. The simulation uses fixed seed \texttt{seed=123} for reproducibility.

\begin{center}
\begin{tabular}{lc}
\hline
Mechanism & Average Rank (lower is better) \\
\hline
DA & 1.2086 \\
IA & 1.1837 \\
TTC & 1.1934 \\
\hline
\end{tabular}
\end{center}

\paragraph{Interpretation}
Across 1000 simulated markets, IA has the lowest average rank (1.1837), so it performs best by this preference-based metric. TTC is close behind at 1.1934 (a gap of 0.0097), indicating similar overall assignment quality in expectation. DA is higher at 1.2086, which is 0.0249 worse than IA on average. Because lower rank means students are assigned closer to their top choices, these differences summarize welfare trade-offs in a compact, reproducible way.

\paragraph{Reproducibility}
From the repository root, run \texttt{python problem\_set\_02/generate\_part1\_3\_writeup.py} to regenerate both Markdown and LaTeX outputs.
 into a larger document.

\section*{Problem Set 02: Parts 1--3}

\subsection*{Part 1: Preferences and Priorities Summary}
The market instance uses 18 students, 3 schools, and school capacity 6 (total capacity 18), with strict random preferences and strict random school priorities. The one-market realization is generated with \texttt{random.seed(42)}. Detailed preference and priority lists are documented in \texttt{problem\_set\_02/market\_setup\_output.md}.

\subsection*{Part 2: One-Market Matching Table (Seed = 42)}
The table below reports student assignments under Deferred Acceptance (DA), Immediate Acceptance (IA/Boston), and Top Trading Cycles (TTC), all run on the same Part 1 market instance.

\begin{center}
\begin{tabular}{lccc}
\hline
Student & DA & IA & TTC \\
\hline
i1 & s2 & s2 & s2 \\
i2 & s3 & s3 & s3 \\
i3 & s3 & s3 & s3 \\
i4 & s3 & s3 & s3 \\
i5 & s1 & s1 & s1 \\
i6 & s1 & s1 & s1 \\
i7 & s3 & s3 & s3 \\
i8 & s2 & s2 & s2 \\
i9 & s2 & s2 & s2 \\
i10 & s2 & s2 & s2 \\
i11 & s1 & s1 & s1 \\
i12 & s3 & s3 & s3 \\
i13 & s1 & s1 & s1 \\
i14 & s2 & s2 & s2 \\
i15 & s1 & s1 & s1 \\
i16 & s1 & s1 & s1 \\
i17 & s2 & s2 & s2 \\
i18 & s3 & s3 & s3 \\
\hline
\end{tabular}
\end{center}

\paragraph{One-Market Average Rank Check}
\begin{center}
\begin{tabular}{lc}
\hline
Mechanism & Average Rank in this one market \\
\hline
DA & 1.2222 \\
IA & 1.2222 \\
TTC & 1.2222 \\
\hline
\end{tabular}
\end{center}

\subsection*{Part 3: Simulation Results (N = 1000)}
Each simulation draw generates a fresh random market, then computes DA/IA/TTC outcomes and student average rank. The simulation uses fixed seed \texttt{seed=123} for reproducibility.

\begin{center}
\begin{tabular}{lc}
\hline
Mechanism & Average Rank (lower is better) \\
\hline
DA & 1.2086 \\
IA & 1.1837 \\
TTC & 1.1934 \\
\hline
\end{tabular}
\end{center}

\paragraph{Interpretation}
Across 1000 simulated markets, IA has the lowest average rank (1.1837), so it performs best by this preference-based metric. TTC is close behind at 1.1934 (a gap of 0.0097), indicating similar overall assignment quality in expectation. DA is higher at 1.2086, which is 0.0249 worse than IA on average. Because lower rank means students are assigned closer to their top choices, these differences summarize welfare trade-offs in a compact, reproducible way.

\paragraph{Reproducibility}
From the repository root, run \texttt{python problem\_set\_02/generate\_part1\_3\_writeup.py} to regenerate both Markdown and LaTeX outputs.
